\begin{solapartat}
Cal que a la superfície superior de la bobina hi hagi un pol SUD \punts{0,3} per tal que atregui el pol NORD de l'imant. El sentit del corrent a la bobina es pot descriure de diferents maneres, ambdues es donaran com a correctes

\begin{minipage}[c]{0.66\linewidth}
\begin{itemize}
  \item[$\checkmark$] Vist des de dalt, el corrent girarà en sentit horari (vist des de sota en sentit antihorari).
\end{itemize}
\end{minipage}\hfill
\begin{minipage}[c]{0.14\linewidth}
\figura[0.7]{sol_fig1.png}
\end{minipage}\hfill
\begin{minipage}[c]{0.14\linewidth}
\punts{0,2}
\end{minipage}

O bé,

\begin{minipage}[c]{0.66\linewidth}
\begin{itemize}
  \item[$\checkmark$] tal com s'indica a la figura (orientada com a l'enunciat):
\end{itemize}
\end{minipage}\hfill
\begin{minipage}[c]{0.3\linewidth}
\figura[0.75]{sol_fig2.png}
\end{minipage}

\begin{minipage}[c]{0.66\linewidth}
\punts{0,5} Les línies del camp magnètic es representen al dibuix
\end{minipage}\hfill
\begin{minipage}[c]{0.3\linewidth}
\figura[0.9]{sol_fig3.png}
\end{minipage}
\end{solapartat}

\begin{solapartat}
\punts{0,4} b1) Augmentaria el camp magnètic produït per la bobina, que atrauria l'imant amb MÉS força.

\punts{0,3} b2) En augmentar la permeabilitat magnètica, augmentaria el camp magnètic produït per la bobina, que atrauria l'imant amb MÉS força.

\punts{0,3} b3) El camp magnètic canviaria de sentit molt ràpidament. L'imant serà successivament atret i repel·lit (potser observaríem una vibració amb la freqüència del corrent altern).
\end{solapartat}
